% Pollis drill-down: OCR Files and Snippets, reading scanned works in other languages: a record for each work, a reader
% that shows each page of the scan beside its text, passages cut out and corrected as snippets, their translation by a
% local model and their reading aloud (25 November 2026). A draft, marked as one on its first page and in its footer,
% until the feature is built; its examples are to be adapted as the feature is built.
% Compile with: latexmk -pdf workflow_ocr_files.tex
% Shared look for the one-page Pollis workflow handouts (same fonts and colours as the brochure).
\documentclass[10pt,a4paper]{article}

\usepackage[a4paper,left=17mm,right=17mm,top=14mm,bottom=16mm,footskip=8mm]{geometry}
\usepackage[T1]{fontenc}
\usepackage[default,semibold]{sourcesanspro}
\usepackage{microtype}
\usepackage{graphicx}
\usepackage[table]{xcolor}
\usepackage{tikz}
\usepackage[most]{tcolorbox}
\usepackage{fontawesome5}
\usepackage{enumitem}
\usepackage{fancyhdr}
\usepackage[hidelinks]{hyperref}

\graphicspath{{images/}}

% ---------------------------------------------------------------- colours (as in the brochure)
\definecolor{navy}{HTML}{1C3A5E}
\definecolor{pblue}{HTML}{3B78B5}
\definecolor{amber}{HTML}{F2B33D}
\definecolor{deepamber}{HTML}{D9902A}
\definecolor{pale}{HTML}{EEF3F9}
\definecolor{ink}{HTML}{2A2F36}
\definecolor{logonavy}{HTML}{173D63}
\definecolor{logodark}{HTML}{C98A1E}
\definecolor{logoamber}{HTML}{F2B446}
\definecolor{logocream}{HTML}{FEEFD2}
\definecolor{soft}{HTML}{5E6773}
\color{ink}

\setlength{\parindent}{0pt}
\setlength{\parskip}{3pt}
\setlist{nosep,leftmargin=*,itemsep=2pt}

\pagestyle{fancy}
\fancyhf{}
\renewcommand{\headrulewidth}{0pt}
\fancyfoot[L]{\footnotesize\color{soft}\textbf{Pollis}\enspace Bright Insights through Exploration, Modelling, Simulation and Optimisation}
\fancyfoot[R]{\footnotesize\color{soft}github.com/saragga/pollis/releases/latest}

% The Pollis logo, redrawn so it prints sharply at any size
\newcommand{\pollislogo}[1]{%
  \begin{tikzpicture}[x=#1/512,y=-#1/512,baseline=(current bounding box.center)]
    \fill[logonavy] (0,0) rectangle (512,512);
    \fill[logodark] (36,422) -- (104,256) -- (172,422) -- cycle;
    \fill[logoamber] (104,422) -- (203,180) -- (302,422) -- cycle;
    \fill[logocream] (203,422) -- (339,89) -- (476,422) -- cycle;
  \end{tikzpicture}}

% Top of the page: logo, name, kind of document (Workflow unless given, e.g. [Drill-Down]) and its subject;
% then the question and a one-line promise
\newcommand{\workflowhead}[4][Workflow]{%
  \thispagestyle{fancy}%
  \begin{minipage}[c]{13mm}\pollislogo{11mm}\end{minipage}%
  \begin{minipage}[c]{\dimexpr\linewidth-13mm\relax}
    {\color{navy}\Large\bfseries Pollis}\enspace{\color{soft}#1}\hfill
    {\color{deepamber}\small\bfseries\MakeUppercase{#2}}
  \end{minipage}\par\vspace{7mm}
  {\color{navy}\fontsize{22}{26}\selectfont\bfseries #3}\par\vspace{4pt}
  {\color{soft}\large #4}\par\vspace{5mm}}

% A menu path: \menu{Explore}{Access Data Libraries}... is written as \path{Explore \gt Access Data Libraries}
\newcommand{\gt}{\,{\color{deepamber}\bfseries ›}\,}
\newcommand{\menupath}[1]{{\footnotesize\bfseries\color{navy}#1}}

% A screenshot if images/<file> exists, otherwise an orange placeholder that says what to capture
\newlength{\shotheight}\setlength{\shotheight}{43mm}
\newcommand{\stepshot}[2]{%
  \IfFileExists{images/#1}{%
    \tcbox[enhanced, colback=black, colframe=black!70, boxrule=0.4pt, arc=2pt,
      left=0pt,right=0pt,top=0pt,bottom=0pt, boxsep=0pt, halign=center]{%
      \includegraphics[width=\dimexpr\linewidth-1pt\relax,height=\shotheight,keepaspectratio]{#1}}}{%
    \begin{tcolorbox}[enhanced, colback=white, colframe=orange!80, boxrule=0.6pt, arc=2pt,
      height=\shotheight, valign=center, halign=center, left=6pt, right=6pt]
      {\color{orange}\Large\faCamera}\par\smallskip
      {\color{orange}\footnotesize\bfseries Screenshot: #2}\par\smallskip
      {\color{orange!80!black}\scriptsize images/#1}
    \end{tcolorbox}}}

% One step: number, title, menu path, image file, what to capture, what it means
\newcommand{\step}[6]{%
  \begin{tcolorbox}
    {\color{deepamber}\Large\bfseries #1}\enspace{\color{navy}\large\bfseries #2}\par\vspace{1pt}
    \menupath{#3}\par\vspace{4pt}
    \stepshot{#4}{#5}\par\vspace{3pt}
    {\small #6}
  \end{tcolorbox}}

\newenvironment{steps}{%
  \begin{tcbraster}[raster columns=2, raster equal height=rows, raster column skip=4mm,
    raster row skip=4mm, enhanced, colback=pale, colframe=pale, boxrule=0pt, sharp corners,
    left=6pt,right=6pt,top=6pt,bottom=6pt, fontupper=\raggedright]}{\end{tcbraster}}

% Boxes for the bottom of the page
\newtcolorbox{pbox}[2][]{enhanced, colback=pale, colframe=amber, boxrule=0pt,
  leftrule=3pt, sharp corners, left=7pt, right=7pt, top=6pt, bottom=6pt,
  fontupper=\small\raggedright, before upper={{\color{navy}\normalsize\bfseries #2}\par\smallskip}, #1}
\newtcolorbox{nbox}[1][]{enhanced, colback=navy, colframe=navy, boxrule=0pt,
  sharp corners, left=8pt, right=8pt, top=6pt, bottom=6pt, coltext=white,
  fontupper=\small\raggedright, #1}

% A Next Steps item inside the bottom boxes
\newcommand{\nsitem}[2]{\item[\color{deepamber}\faChevronRight] \textbf{#1}: #2}

% Citations: \dcite{key}{Breiman (2001)} links to its entry \dref{key}{...} in the References, keyed as in pollis.bib
\newcommand{\dcite}[2]{\hyperlink{ref:#1}{\color{pblue}#2}}
\newcommand{\dref}[2]{\raisebox{1.1em}[0pt][0pt]{\hypertarget{ref:#1}{}}#2\par}
\newcommand{\ddoi}[1]{\href{https://doi.org/#1}{\color{pblue}doi:#1}}
\newcommand{\references}{\section*{References}\addcontentsline{toc}{section}{References}\footnotesize\raggedright\setlength{\parskip}{3pt}}

% The page of each heading, for Ask Pollis
% Pollis: the page of each heading, for Ask Pollis. A document that inputs this file, at the end of its preamble, writes
% <jobname>.labels beside its PDF, a line page|level|heading for each heading, the page counted from the first, the
% cover, as 1. compile_documents.jl copies it to ~/.pollis/documents, and ask_index.jl cuts the document into passages
% under its headings with it. \asklabel{1}{Title} marks a topic or a page, \asklabel{2}{Title} a heading inside it; the
% sections and subsections of a document that has them are marked as they go into its contents. A package vignette
% inputs it before its \documentclass (compile_vignette.jl), so the contents are taken over when the document begins,
% after the packages have taken them over. Read a second time, as by a copy of a document for Edit the LaTeX Source,
% which inputs it as resources/ask_labels when Pollis has put it in already (document_source.jl), it does nothing.
\ifdefined\asklabel\endinput\fi
\makeatletter
\newwrite\pollis@labels
\immediate\openout\pollis@labels=\jobname.labels
\newcommand{\asklabel}[2]{\write\pollis@labels{\the\ReadonlyShipoutCounter|#1|\unexpanded{#2}}}
\def\pollis@toc{toc}
\newcommand{\pollis@contentsline}[3]{%
  \def\pollis@file{#1}%
  \ifx\pollis@file\pollis@toc\asklabel{#2}{#3}\fi
  \pollis@addcontentsline{#1}{#2}{#3}}
% The headings of the New Methods style of the vignettes, \nmsection and \nmsub, which go into no contents
\def\pollis@asksection#1{\pollis@nmsection{#1}\asklabel{1}{#1}}
\def\pollis@asksub#1{\pollis@nmsub{#1}\asklabel{2}{#1}}
\AddToHook{begindocument/end}{%
  \let\pollis@addcontentsline\addcontentsline\let\addcontentsline\pollis@contentsline
  \ifdefined\nmsection\let\pollis@nmsection\nmsection\let\nmsection\pollis@asksection\fi
  \ifdefined\nmsub\let\pollis@nmsub\nmsub\let\nmsub\pollis@asksub\fi}
\makeatother



\usepackage{titlesec}
\usepackage{tabularx}
\usepackage{booktabs}
\usepackage{listings}
\usepackage{needspace}

\fancyhf{}
\fancyfoot[L]{\footnotesize\color{soft}\textbf{Pollis}\enspace Drill-down: OCR Files and Snippets\enspace·\enspace Draft}
\fancyfoot[R]{\footnotesize\color{soft}\thepage}

\titleformat{\section}{\color{navy}\large\bfseries}{\color{deepamber}\thesection}{0.6em}{}
\titleformat{\subsection}{\color{navy}\normalsize\bfseries}{\color{deepamber}\thesubsection}{0.5em}{}
\titlespacing*{\section}{0pt}{12pt}{4pt}
\titlespacing*{\subsection}{0pt}{8pt}{2pt}
\setlength{\parskip}{4pt}

\lstset{basicstyle=\ttfamily\footnotesize, columns=fullflexible, keepspaces=true, breaklines=true,
  backgroundcolor=\color{pale}, frame=leftline, rulecolor=\color{amber}, framerule=2pt, xleftmargin=6pt,
  framexleftmargin=6pt, aboveskip=4pt, belowskip=4pt, upquote=true, tabsize=2,
  literate={é}{{\'e}}1 {è}{{\`e}}1 {à}{{\`a}}1 {ç}{{\c{c}}}1 {£}{{\pounds}}1 {—}{{---}}1}

\usepackage{xurl}
\newcommand{\file}[1]{{\small\nolinkurl{#1}}}
\newcolumntype{L}{>{\raggedright\arraybackslash}X}
\newcommand{\ui}[1]{\textbf{#1}}

\begin{document}

\workflowhead[Drill-Down\enspace{\color{deepamber}\bfseries Draft}]{OCR Files and Snippets}
  {Reading scanned works in other languages}
  {A record for each work, its pages beside their text, passages corrected as snippets, translated by a local model
   and read aloud.}

\begin{nbox}
{\color{amber}\bfseries Drill-down, draft.}\enspace
This drill-down describes OCR Files and Snippets before the feature is in Pollis: its menu items are not there yet,
and its examples will change with it.

Many sources are old printed books in other languages, read from scans whose text has errors. OCR Files and Snippets,
beside PDF Files, keeps a record for each scanned work, shows each page of the scan beside its text, cuts passages out
and corrects them as snippets, and translates them with a local model you choose, without the internet.

This drill-down describes each part: the work's record (\S\ref{sec:record}), its contents (\S\ref{sec:contents}), the
checks that find misread words (\S\ref{sec:correct}), snippets (\S\ref{sec:snippets}), translation
(\S\ref{sec:translation}), reading aloud (\S\ref{sec:aloud}), where works are kept (\S\ref{sec:where}) and the menu
items (\S\ref{sec:menu}).
\end{nbox}

\tableofcontents

\section{What it is for}\label{sec:why}
\begin{itemize}
  \item \textbf{The original language is often the only free text.} For Van Gogh's last years, 1886 to 1890, the scan
    with its text is volume 3 of \emph{Brieven aan zijn broeder} (1914). Its letters are mostly in French, with some
    in Dutch.
  \item \textbf{The scan's text has errors that change meaning}: ``s6rie'' for s\'erie, ``\pounds tude'' for
    \'etude, ``Yerde'' for Verde, ``Auversest'' for Auvers est.
  \item \textbf{Errors in the text become errors in the translation.} A model translates well only text that has been
    read right, so a passage is corrected before it is translated.
\end{itemize}
The feature is not tied to one toolbox: archives, old editions, inscriptions and the Britannica page have the same
need.

\section{How it works}\label{sec:how}
The work's record says what the work is and where its scan is; every step reads it.
\begin{enumerate}
  \item \textbf{Open a work}: from its Internet Archive identifier, a Project Gutenberg number or a local PDF. Pollis
    writes its record, filling what the source gives: title, author, year, language, contents.
  \item \textbf{Read it}: each page of the scan beside its text, with the words the checks flag highlighted.
  \item \textbf{Choose parts}: one or more entries of the contents, or a page range.
  \item \textbf{Extract and correct}: the text of the chosen parts, with each flagged word beside its piece of the
    scan, corrected by you. The result is a snippet.
  \item \textbf{Translate}: the snippet, with a local model you choose, into the language you pick. The translation
    is saved with the snippet, labelled with its model and date.
  \item \textbf{Read aloud}: the snippet, or its translation.
  \item \textbf{Use it}: copy a snippet, send it to a notebook, or cite it from a panel.
\end{enumerate}

\section{The work record}\label{sec:record}
Each work is one TOML file, \file{work.toml}, in its own folder \file{~/.pollis/works/<id>/}, beside its snippets.
Every field but \file{id} and the scan is optional, since some works have no known author, year or even title.

\begin{center}\small
\begin{tabularx}{\linewidth}{@{}p{0.29\linewidth}L@{}}
\toprule
\textbf{Field} & \textbf{What it holds} \\
\midrule
\file{id} & A short name for the folder, such as \file{van-gogh-brieven-3} \\
\file{title}, \file{author}, \file{editor}, \file{translator}, \file{year} & As printed, when known \\
\file{languages} & The languages of the text, as codes (\file{fr}, \file{nl}), the language of most of it first \\
\file{entryType} & The BibTeX entry type: \file{book}, \file{article}, \file{incollection}, \file{misc}, and so on \\
\file{referenceKey} & The work's key in PollisReferences, when it is there; the record then does not repeat its
  details \\
\file{summary}, \file{summarySource} & A few sentences on the work, and the page they come from, such as a Wikipedia
  article or a catalogue \\
\file{[scan]} & Where the scan and its text are: an Internet Archive identifier, a Gutenberg number or a local file
  (\S\ref{sec:where}) \\
\file{[[contents]]} & The table of contents, each entry with its label and pages (\S\ref{sec:contents}) \\
\file{[[words]]} & The work's own words that the dictionary does not know, accepted once: names, old spellings \\
\file{[[glossary]]} & Agreed translations of names and recurring terms, used by every translation of the work \\
\bottomrule
\end{tabularx}
\end{center}

The record of the third volume of Van Gogh's letters to Theo:
\begin{lstlisting}
id = "van-gogh-brieven-3"
title = "Brieven aan zijn broeder, derde deel"
author = "Vincent van Gogh"
editor = "J. van Gogh-Bonger"
year = 1914
languages = ["fr", "nl"]
entryType = "book"
summary = "The letters of Vincent van Gogh to his brother Theo from Paris, Arles, Saint-Rémy and Auvers, 1886 to 1890, as he wrote them, with his sister-in-law's notes in Dutch."

[scan]
archive = "brievenaanzijnbr03gogh"

[[contents]]
label = "Verblijf in Parijs. Maart 1886 — Februari 1888"
pages = "1-5"

[[contents]]
label = "Verblijf in Arles. Februari 1888 — Mei 1889"
pages = "6-"

[[words]]
word = "Gachet"

[[glossary]]
term = "le midi"
en = "the South"
\end{lstlisting}

\section{Contents}\label{sec:contents}
Each contents entry points to pages of the scan, so choosing parts works the same for every work. You pick one entry or
several.

\subsection{Where the contents come from}
\begin{enumerate}
  \item \textbf{The Internet Archive's contents file} (\file{_toc.xml}), the OCR of the book's own contents page. It
    is a draft to check: for Van Gogh's volume 3 it reads ``Verblijf in Arles. Februari 1888 --- Mei 1889. Brief
    463---590 - 6'', but some of its lines run a page number into the next entry.
  \item \textbf{The PDF's outline}, for a local PDF that has one.
  \item \textbf{The headings} of a Gutenberg text.
  \item \textbf{You}, entry by entry, in the reader.
\end{enumerate}

\subsection{Without a table of contents}
The entries are other units of the work: Van Gogh's numbered letters, the sections of the Periplus of the Erythraean
Sea, the stanzas of the Meghaduta. When nothing else is possible, the entries are page ranges.

\subsection{Printed pages and scan pages}
They differ: a scan's first pages are covers and front matter, and plates add pages with no number. The Internet
Archive's page-numbers file (\file{_page_numbers.json}) gives the printed number of each scan page, with the
probability of its reading: in volume 3, scan page 32 is printed page 20, read with a probability of 93 per cent.
Contents entries give printed pages, and Pollis turns them into scan pages. Where there is no such file, you set the
difference once for the work.

\section{Extracting and correcting}\label{sec:correct}
The checks flag the words most likely misread; you correct them, each beside its piece of the scan. Nothing is changed
without you.

\begin{center}\small
\begin{tabularx}{\linewidth}{@{}p{0.2\linewidth}LL@{}}
\toprule
\textbf{Check} & \textbf{What it catches} & \textbf{Example} \\
\midrule
The OCR's own confidence & Words the OCR itself was unsure of: the Internet Archive's hOCR file gives every word its
  box on the page and a confidence from 0 to 100 & ``ziekenkiiizen'' (ziekenhuizen), confidence 7 \\
A dictionary of the language & Words in no dictionary, from Hunspell's free word lists, which LibreOffice uses
  & ``s6rie'', ``Auversest'' \\
Word lists of the work's time & Old spellings that a modern dictionary flags, where such lists exist: Portuguese
  spelling, for instance, changed in 1911, 1945 and 1990 & ``gelukwenschen'' in Dutch before 1947 \\
The work's own words & Names and spellings accepted once, so they are not flagged again & Gachet, Roulin \\
Common OCR confusions & The correction of a flagged word: \'e read as 6, \`e as \pounds, c as e, rn as m, li as h &
  ``eher'' $\rightarrow$ cher, ``\pounds tude'' $\rightarrow$ \'etude \\
A second scan & Misreads that are real words, where two scans of the same edition disagree & A word one scan reads
  right \\
The language model & Words that do not fit their sentence, as suggestions only & A misread that is still a word \\
\bottomrule
\end{tabularx}
\end{center}

Every correction is kept as ``read as X, corrected to Y'', so the scan's reading can always be restored. A second scan
is the strongest check, and the Internet Archive often has more than one copy of a book: Pollis offers the other copies
of the same edition it finds.

\subsection{Beyond printed text}
The checks above are for printed text in a modern script, whose scan already has its text. A page with no text, such
as a page you photograph, needs OCR first. An inscription in an ancient script is a case of its own: its extraction
and translation need an OCR or a vision model that reads that script.

\needspace{16\baselineskip}
\section{Snippets}\label{sec:snippets}
A snippet is a passage with its place in the work, its corrected text, the corrections made and any translations, of
any size you choose: a couple of paragraphs, a chapter, one inscription. It is one TOML file in the work's
\file{snippets} folder:
\begin{lstlisting}
work = "van-gogh-brieven-3"
label = "Brief 499, Saintes-Maries"
letter = 499
language = "fr"
text = "Je t'écris de Stes. Maries au bord de la Méditerranée enfin."

[[corrections]]
read = "M£diterranee"
corrected = "Méditerranée"

[[translations]]
language = "en"
model = "EuroLLM-9B-Instruct-2512.Q4_K_M"
made = 2026-11-25
text = "I write to you from Stes. Maries on the Mediterranean at last."
\end{lstlisting}
The passages of the Historical Geography topics are snippets in all but name: a work, a place in it and its words. A
topic's passage becomes a snippet when you make it one.

\section{Translation}\label{sec:translation}
A snippet is translated by a local model you choose, run by llama-server as Pollis runs its other local models, so it
works without the internet. The translation is saved with the snippet, so each passage is translated once.

\subsection{Context without losing time}
Writing a translation is slow, about 6 words a second for a model of 9 billion parameters on a Mac with 16 GB of
memory; reading the source is much faster. So each request carries more than it translates:
\begin{itemize}
  \item the paragraph before, with its translation already made;
  \item the paragraph to translate;
  \item the paragraph after, marked as context only;
  \item the work's glossary.
\end{itemize}
Short paragraphs, as in letters and poems, are grouped into units of a few hundred words. The translation is shown word
by word as it is written. \ui{Translate Again} translates a passage with the whole snippet as context, for the passage
that only later paragraphs explain.

\subsection{The model}
The model must take instructions: TranslateGemma's fixed prompt has no room for context or a glossary, while general
models such as EuroLLM and Gemma~4 take both. Three models on four passages of the 1914 letters, three in French and one
in Dutch, the scan's errors corrected first:

\begin{center}\small
\begin{tabularx}{\linewidth}{@{}p{0.22\linewidth}LLL@{}}
\toprule
\textbf{Passage} & \textbf{TranslateGemma 4B} & \textbf{Gemma 4 12B} & \textbf{EuroLLM 9B-2512} \\
\midrule
Arles, on arrival (French) & Right & Right & Right \\
Saintes-Maries (French) & ``sardines'' for maquereaux (mackerel) & Right, one clumsy phrase & Right \\
Auvers, the first letter (French) & ``old reeds'' for vieux chaume (old thatch) & Right & Right \\
To his mother (Dutch) & Reverses the meaning of a sentence & Gives his mother's move to Vincent & The right person,
  one nuance off \\
\midrule
Seconds a passage & 9 to 12 & 28 to 51 & 22 to 28 \\
Size at 4 bits & 2.5 GB & 7.1 GB & 5.6 GB \\
\bottomrule
\end{tabularx}
\end{center}

The times were measured with a download running and two embedding servers of Pollis in memory. None of the three was
right in every sentence of the Dutch letter, so every translation is labelled as made by a model, with the model and
the date.

\section{Reading aloud}\label{sec:aloud}
\ui{Read Aloud} reads a snippet, or its translation, with the voice tools of Pollis, in the language of the text.

\needspace{8\baselineskip}
\section{Where works are kept}\label{sec:where}
\subsection{One folder per work}
A work's folder, \file{~/.pollis/works/<id>/}, holds two kinds of files:
\begin{itemize}
  \item \textbf{Small text files}, a few kilobytes each: the record, the snippets with their corrections and
    translations, the glossary and the word list. They are what you make.
  \item \textbf{Large files}, kept only as a copy of the source: the scan's text, the word positions (the hOCR file of
    volume 3 is 28 MB) and the page images (221 MB for volume 3). Pollis can fetch each page's image when it is shown,
    or download the whole scan to read offline; it shows the size of the scan and the free disk space, and you choose.
\end{itemize}

\subsection{A work in a project}
\ui{Add to Workspace} puts a work into the Pollis Workspace open, so that a project keeps its sources when it is copied
or shared. It copies only the small text files. The large files stay in \file{~/.pollis/works/}, once on the computer,
and the record points to them by a permalink, the address that names the source for good:
\begin{center}\small
\begin{tabularx}{\linewidth}{@{}p{0.2\linewidth}p{0.36\linewidth}L@{}}
\toprule
\textbf{Source} & \textbf{Permalink} & \textbf{On another computer} \\
\midrule
Internet Archive & The item's identifier: \file{archive = "brievenaanzijnbr03gogh"}, that is,
  \file{archive.org/details/brievenaanzijnbr03gogh} & Pollis fetches the text and the pages again from the same address \\
Project Gutenberg & The ebook's number: \file{gutenberg = 16659} & The same \\
A local PDF & Its fingerprint, the SHA-256 of the file, with the path it had & Pollis finds the file in
  \file{~/.pollis/works/} by its fingerprint, or asks for it \\
\bottomrule
\end{tabularx}
\end{center}
Corrections and translations belong to the snippets, not to the large files, so nothing you made is lost when the large
files are fetched again.

\section{Menu items}\label{sec:menu}
The submenu \ui{Explore}\gt\ui{OCR Files and Snippets}, below \ui{PDF Files}. Its page has numbered steps, those of
\S\ref{sec:how}; the other items open it at a step or work on what is open.

\begin{center}\small
\begin{tabularx}{\linewidth}{@{}p{0.22\linewidth}L@{}}
\toprule
\textbf{Item} & \textbf{What it does} \\
\midrule
\ui{Open Work...} & Asks for an Internet Archive identifier or link, a Gutenberg number or a local PDF, writes the
  work's record and opens the page at Read \\
\ui{Works} & Lists the works opened, with their title, author, year and languages \\
\ui{Read} & The scan's pages beside their text, the flagged words highlighted, the contents at the side \\
\ui{Extract Snippet} & The chosen contents entries or pages as a snippet, at the step of correcting it \\
\ui{Translate Snippet} & The snippet translated, with the model and the language you choose \\
\ui{Read Aloud} & Reads the snippet, or its translation, aloud \\
\ui{Snippets} & Lists the snippets of every work, with their translations, to copy, send to a notebook or open \\
\ui{Glossary} & The work's agreed translations of names and terms, to edit \\
\ui{Add to Workspace} & Puts the work's record and snippets into the Pollis Workspace open (\S\ref{sec:where}) \\
\bottomrule
\end{tabularx}
\end{center}

\section*{Sources}\addcontentsline{toc}{section}{Sources}
\begin{itemize}
  \item \href{https://archive.org/details/brievenaanzijnbr03gogh}{\color{pblue}Brieven aan zijn broeder, volume 3
    (1914)}, the Getty Research Institute's copy at the Internet Archive, with its
    \href{https://archive.org/metadata/brievenaanzijnbr03gogh/files}{\color{pblue}files}: text, hOCR, contents and
    page numbers.
  \item \href{https://huggingface.co/utter-project/EuroLLM-9B-Instruct-2512}{\color{pblue}EuroLLM-9B-Instruct-2512},
    and its \href{https://huggingface.co/mradermacher/EuroLLM-9B-Instruct-2512-GGUF}{\color{pblue}4-bit GGUF
    conversion}.
\end{itemize}

\end{document}
